\chapter{Multi-Criteria Decision Analysis for DSS}
\label{ch:multi-criteria-decision-analysis-for-dss}

\section{Chapter Overview and Learning Objectives}

Many important DSS problems cannot be reduced to a single objective. Managers
often balance cost, quality, risk, fairness, time, and strategic fit at the
same time. Multi-Criteria Decision Analysis (MCDA) provides structured methods
for evaluating alternatives when several criteria matter simultaneously.

\subsection*{Learning objectives}
\addcontentsline{toc}{subsection}{Learning objectives}
By the end of this chapter, the student should be able to:
\begin{itemize}
  \item explain why MCDA is important in semi-structured DSS problems,
  \item construct a criteria matrix and normalize decision criteria,
  \item compare weighted scoring, AHP, and outranking-style reasoning at a high
  level,
  \item perform sensitivity analysis on criteria weights,
  \item identify governance risks in criteria design and weighting.
\end{itemize}

\section{Why Multi-Criteria Methods Matter}

\subsection{Single-objective models are often insufficient}
Real DSS problems often involve competing objectives:
\begin{itemize}
  \item selecting a supplier with trade-offs among cost, quality, and delivery,
  \item prioritizing projects under risk, benefit, and strategic alignment,
  \item ranking interventions where effectiveness, equity, and feasibility all
  matter.
\end{itemize}

\subsection{Core MCDA workflow}
A practical MCDA workflow in DSS includes:
\begin{enumerate}[label=\arabic*.]
  \item define alternatives,
  \item define criteria and measurement scales,
  \item determine directions of preference,
  \item assign or elicit weights,
  \item aggregate or compare alternatives,
  \item test sensitivity and robustness,
  \item document rationale and final decision.
\end{enumerate}

\section{Weighted Scoring Models}

\subsection{Decision matrix}
The starting point is a decision matrix where rows represent alternatives and
columns represent criteria. Because criteria can have different units,
normalization is often required before aggregation.

\subsection{Weighted sum model}
One of the simplest approaches is the weighted sum:
$$
S(a) = \sum_{j=1}^{m} w_j x_{aj},
$$
where $w_j$ is the weight of criterion $j$ and $x_{aj}$ is the normalized score
of alternative $a$ on criterion $j$.

\subsection{Strengths and limitations}
Weighted scoring is transparent and easy to implement in DSS dashboards, but it
assumes compensability: poor performance on one criterion can be offset by high
performance on another. In some settings, this is undesirable.

\section{Analytic Hierarchy Process (AHP)}

\subsection{Hierarchical structuring}
AHP decomposes the problem into goal, criteria, subcriteria, and alternatives.
Decision makers compare criteria pairwise, which can be more intuitive than
assigning direct numeric weights.

\subsection{Consistency}
An advantage of AHP is that it checks the consistency of pairwise judgments. If
judgments are highly inconsistent, the DSS can prompt the user to revise them
before ranking alternatives.

\section{Sensitivity and Robustness}

\subsection{Why sensitivity analysis is essential}
Small changes in weights can reverse rankings. A DSS should therefore expose:
\begin{itemize}
  \item how rankings change when weights change,
  \item which criteria dominate the result,
  \item where alternatives are nearly tied.
\end{itemize}

\subsection{Scenario-based MCDA}
One useful pattern is to define stakeholder-specific or scenario-specific
weighting schemes. For example, a finance perspective may weight cost more
heavily, while a service perspective may weight quality and responsiveness more
heavily.

\section{Governance Considerations}

\subsection{Criteria design is a policy decision}
MCDA can appear objective while embedding subjective choices. Governance
therefore requires transparent documentation of:
\begin{itemize}
  \item why criteria were selected,
  \item how weights were chosen,
  \item who approved the weighting scheme,
  \item how conflicts are resolved.
\end{itemize}

\subsection{Fairness and accountability}
If criteria disadvantage specific groups, the DSS may institutionalize bias.
Criteria and weights should therefore be reviewed for fairness, legality, and
organizational legitimacy.

\section{Summary and Concluding Remarks}

MCDA provides a disciplined way to evaluate alternatives when several criteria
must be balanced. In DSS practice, its value is not only computational but also
communicative: it makes trade-offs explicit, supports stakeholder discussion,
and documents why an alternative was ranked highly. The most important design
principle is transparency in criteria, weights, and sensitivity analysis.

\section{Key Terms}
\begin{description}[style=nextline]
  \item[MCDA] Multi-Criteria Decision Analysis; a family of methods for
  evaluating alternatives across several criteria.
  \item[Decision matrix] Table of alternatives and their scores on criteria.
  \item[Normalization] Transforming criteria values to comparable scales.
  \item[Weight] Relative importance assigned to a criterion.
  \item[Compensability] The extent to which poor performance on one criterion
  can be offset by strong performance on another.
  \item[AHP] Analytic Hierarchy Process; a structured MCDA method using pairwise
  comparisons.
  \item[Sensitivity analysis] Studying how changes in assumptions affect the
  ranking or recommendation.
\end{description}

\section{Multiple-Choice Questions (MCQs)}

\subsection*{Questions}
\addcontentsline{toc}{subsection}{Questions}
\begin{enumerate}[label=\textbf{Q\arabic*.}]
  \item MCDA is especially useful when:
  \begin{enumerate}[label=\alph*)]
    \item there is only one criterion and no constraints
    \item decision makers must balance multiple competing criteria
    \item labels are unavailable for clustering
    \item a rule base replaces all judgment
  \end{enumerate}

  \item A weighted scoring model requires:
  \begin{enumerate}[label=\alph*)]
    \item criteria, scores, and weights
    \item only a confusion matrix
    \item only symbolic rules
    \item no normalization at all
  \end{enumerate}

  \item Sensitivity analysis helps determine:
  \begin{enumerate}[label=\alph*)]
    \item whether rankings are robust to weight changes
    \item whether a database is normalized
    \item whether an image classifier is deep enough
    \item whether all alternatives have equal cost
  \end{enumerate}

  \item AHP is known for:
  \begin{enumerate}[label=\alph*)]
    \item pairwise comparison of criteria and consistency checking
    \item only using random selection
    \item ignoring stakeholder judgment
    \item replacing all decision matrices
  \end{enumerate}
\end{enumerate}

\subsection*{Answer Key}
\addcontentsline{toc}{subsection}{Answer Key}
\begin{enumerate}[label=\textbf{Q\arabic*.}]
  \item b
  \item a
  \item a
  \item a
\end{enumerate}
